% TOP_PROBLEM: {"id": 30007499, "title": "Superpolynomial worst-case trace sample lower bounds", "category_id": 15, "category": {"id": 15, "name": "computer_science", "display_name": "Computer Science", "description": "Computational complexity, algorithms, and theoretical CS.", "slug": "computer-science", "order_index": 15, "created_at": "2026-07-31T15:26:25.671Z"}, "status": "open"}
% FIRST_POSTED: 2026-10-06
% ATTEMPT_STATUS: unresolved
% OPENAI_CLAIM: {"schema_version":1,"source_repo":"https://github.com/openai/math","source_commit":"adc7f1241b42e322a6451854ab7e4b4c146bf78a","release_date":"2026-10-06","match":"partial","resolution":"partial","independently_reviewed":false,"second_pass":"n/a","families":[{"family":"122","title":"Quantitative trace-reconstruction bounds with a uniform decoder","subject":"Theoretical computer science","manuscripts":[{"title":"Quantitative lower bounds for trace reconstruction","pdf_path":"preprints/quantitative-lower-bounds-for-trace-reconstruction-September-24-2026/paper.pdf","readme_path":"preprints/quantitative-lower-bounds-for-trace-reconstruction-September-24-2026/README.md","date":"2026-09-24","theorem_ref":"preprints/quantitative-lower-bounds-for-trace-reconstruction-September-24-2026/build/sections/introduction.tex#thm:main"}],"lean":{"doc_path":"lean/docs/122.md","comparators":["lean/ComparatorChallenges/TraceReconstruction.json","lean/ComparatorChallenges/TraceReconstruction.lean"],"declarations":[],"covers_main_theorem":true},"reasoning_trace":null}]}
\documentclass[11pt]{article}
\usepackage[margin=25mm]{geometry}
\usepackage{fontspec}
\setmainfont{Latin Modern Roman}
\usepackage{amsmath,amssymb,mathtools}
\usepackage{unicode-math}
\setmathfont{Latin Modern Math}
\usepackage{enumitem}
\usepackage{xurl,hyperref,bookmark}
\hypersetup{colorlinks=true,urlcolor=blue}
\setcounter{secnumdepth}{0}
\setlength{\parindent}{0pt}
\setlength{\parskip}{5pt}
\setlength{\emergencystretch}{3em}
\begin{document}
\section*{Superpolynomial worst-case trace sample lower bounds}

\subsection{OpenAI's claim}
OpenAI presents the following claims relevant to this question, with manuscript numbers following paper order in Sources. In manuscript 1 (family 122; 2026-09-24), the asserted result is that worst-case exact binary trace reconstruction needs $n^{\Omega(\log\log n)}$ samples at any fixed deletion probability, even with unrestricted computation.

\subsection{Outline of the argument}
For manuscript 1 (family 122; 2026-09-24), nested signed words hide fair orientation bits inside overlapping windows. An auxiliary observation retains side data while projecting to genuine traces; likelihood-score induction forces a large binary subtree for every surviving term and bounds the distance between the two trace laws.

\subsection{Scope relative to this problem}
The lower bound concerns exact worst-case reconstruction at each fixed deletion probability and positive success probability, allowing unlimited computation. Complete coverage of the precise statistical statement is not established by this catalogue, so it remains partial.

\subsection{Formal verification}
For manuscript 1 (family 122; 2026-09-24), OpenAI lists Lean statements for this assertion: worst-case exact binary trace reconstruction needs $n^{\Omega(\log\log n)}$ samples at any fixed deletion probability, even with unrestricted computation.

The selected statement includes the clause used for this problem.

The linked family Comparator checks are \texttt{TraceReconstruction.lean}.

Their named declarations include \texttt{OAI.TraceReconstruction.quantitative\_sample\_lower\_bound}, \texttt{OAI.TraceReconstruction.superpolynomial\_one\_trace}, \texttt{OAI.TraceReconstruction.superpolynomial\_sample\_complexity}.

This site has not executed the Lean code. Described statement coverage is not an independent proof check.

\subsection{Status caveat}
The manuscripts were produced by an internal OpenAI model. OpenAI cautions that results without formalization can contain mistakes. This site has not independently checked the proofs; the label records the scope of OpenAI's claim.

\subsection{Sources}
\begin{itemize}
\item OpenAI, \emph{Quantitative lower bounds for trace reconstruction}, 2026-09-24; citation key \texttt{OAI:quantitative-lower-bounds-for-trace-reconstruction-September-24-2026}. \href{https://github.com/openai/math/blob/adc7f1241b42e322a6451854ab7e4b4c146bf78a/preprints/quantitative-lower-bounds-for-trace-reconstruction-September-24-2026/paper.pdf}{PDF (GitHub)}; \href{https://github.com/openai/math/raw/adc7f1241b42e322a6451854ab7e4b4c146bf78a/preprints/quantitative-lower-bounds-for-trace-reconstruction-September-24-2026/paper.pdf}{PDF (direct)}; \href{https://github.com/openai/math/blob/adc7f1241b42e322a6451854ab7e4b4c146bf78a/preprints/quantitative-lower-bounds-for-trace-reconstruction-September-24-2026/README.md}{Manuscript page and citation}; \href{https://github.com/openai/math/blob/main/preprints/quantitative-lower-bounds-for-trace-reconstruction-September-24-2026/paper.pdf}{Latest version}.
\item Family 122: \href{https://github.com/openai/math/blob/adc7f1241b42e322a6451854ab7e4b4c146bf78a/lean/docs/122.md}{Lean scope}.
\item \href{https://github.com/openai/math/blob/adc7f1241b42e322a6451854ab7e4b4c146bf78a/lean/ComparatorChallenges/TraceReconstruction.json}{Comparator TraceReconstruction.json}.
\item \href{https://github.com/openai/math/blob/adc7f1241b42e322a6451854ab7e4b4c146bf78a/lean/ComparatorChallenges/TraceReconstruction.lean}{Comparator TraceReconstruction.lean}.
\item \href{https://github.com/openai/math/blob/adc7f1241b42e322a6451854ab7e4b4c146bf78a/CONTENTS.md}{OpenAI catalogue entry}.
\end{itemize}
\end{document}
